






Uhlmann's theorem \cite{uhlmann1976transition} has been playing a crucial role in quantum information
science. The theorem states that for quantum states $\rho^\sA$ and $\sigma^\sA$ on system $\sA$ and their purifications $\ket{\rho}^{\sA\sB}$ and $\ket{\sigma}^{\sA\sB}$, it holds that 
\begin{align}
\label{inteq:57}
    \rF(\rho^\sA, \sigma^\sA) = \max_{U^\sB}\rF\big((\bI^\sA \otimes U^\sB) \ketbra{\rho}{\rho}^{\sA\sB}(\bI^\sA \otimes U^\sB)^\dag, \ketbra{\sigma}{\sigma}^{\sA\sB}\big),
\end{align}
where the maximization is taken over all unitaries acting on the purifying system $\sB$, and $\bI^\sA$ is the identity operator on $\sA$.
The fidelity $\rF(\rho, \sigma) = \big(\tr\big[\sqrt{\sqrt{\sigma}\rho\sqrt{\sigma}}\big]\big)^2$ is a closeness measure of quantum states $\rho$ and $\sigma$; it increases as the states become closer. 
Notably, the Uhlmann's theorem characterizes optimal local transformation between $\ket{\rho}^{\sA\sB}$ and $\ket{\sigma}^{\sA\sB}$ because, for any quantum channel $\cT^\sB$ acting on $\sB$, $\rF\big((\rid^\sA \otimes \cT^\sB)(\ketbra{\rho}{\rho}^{\sA\sB}), \ketbra{\sigma}{\sigma}^{\sA\sB}\big)$ never exceeds $\rF(\rho^\sA, \sigma^\sA)$, where $\rid^\sA$ is the identity map on $\sA$. 
This implies that the unitary achieving the maximization in Eq.~\eqref{inteq:57} is the optimal local transformation on $\sB$, which takes $\ket{\rho}^{\sA\sB}$ as close as theoretically possible to $\ket{\sigma}^{\sA\sB}$. 
This unitary $U^\sB$ is called the Uhlmann unitary, and the resulting transformation is the Uhlmann transformation. The Uhlmann unitary $U^\sB$ is not determined solely by the states $\rho^\sA$ and $\sigma^\sA$, but only once their purifications are specified.
The Uhlmann transformation has long served as a key tool in various fields, including the quantum Shannon theory \cite{devetak2005distillation, Abeyesinghe2009Motherfamilytree, Datta2011apexfamilytree, Datta2013OneShotCompression, bennett2014quantum, Wilde2017ConvrersePrivate, Metger2024entropyaccumulation, mazzola2025uhlmannstheoremrelativeentropies, fang2025variexpressuhlmann}, complexity theory \cite{kitaev2000parallelamplif, arnonfriedman2023computationalentanglementtheory, metger2023stateqipspace, bostanci2023unicompuhlmann, chia2024complexitytheoryquantumpromise, poremba2024LearningStabilizer}, quantum cryptography \cite{Mayer1997unconditionallysecure, Lo1998WhyBitCommit, Yan2022BitCommit, Khurana2024commitment}, and also fundamental physics \cite{hayden2007black, harlow2013CompvsFirewall, aaronson2016moneyblackhole, Kirklin2020holographicdual, Kirklin2022islandUhlmannphase, brakerski2023blackholeradiation}.


Despite its broad applicability, however, the Uhlmann transformation has an underexplored aspect: the Uhlmann's theorem does not offer an explicit way to realize the transformation in the form of quantum circuits, leaving its algorithmic properties open.
In particular, when we do not have prior knowledge of the descriptions of the states $\ket{\rho}^{\sA\sB}$ and $\ket{\sigma}^{\sA\sB}$, yet have multiple access to samples of these states or queries to the unitaries that prepare them, the following question naturally arises:
\begin{center}
\emph{How can we construct a quantum circuit implementing the Uhlmann transformation using these resources?}
\end{center}
As quantum technologies advance and practical interest grows, the importance of these algorithmic features, such as the computational cost, is increasing.



Recently, there has been progress along this line.
In Ref. \cite{metger2023stateqipspace}, a quantum query algorithm was proposed with an explicit quantum circuit for implementing the Uhlmann transformation. Their algorithm realizes the transformation with polynomial-space computation, marking substantial progress. 
However, it incurs unavoidable exponential costs other than space complexity, because it relies on an exponentially large number of purified state preparations and measurements, as well as an exponentially deep quantum circuit (see \Cref{sec:metger and yuen's algorithm} for details).
These high computational costs motivate us to develop more efficient algorithms.


In this paper, we resolve the problem by proposing quantum query and sample algorithms that implement the Uhlmann transformation in the form of a quantum circuit.
Our algorithms are constructed in three commonly studied computational models: the purified query, the purified sample, and the mixed sample access models.
We then analyze their computational costs, including the query and sample complexities, the number of one- and two-qubit gates, and the number of qubits at any one time.
It is shown that our algorithms can significantly outperform the approach in Ref. \cite{metger2023stateqipspace} and other naive approaches based on quantum state tomography.
In particular, when either $\rho^\sA$ or $\sigma^\sA$ is a polynomial-rank state, our algorithms can be implemented with polynomial computational cost up to inverse-polynomial accuracy in purified query and sample access models.
Thus, they achieve an \emph{exponential} improvement.
While previous approaches rely on a large number of state measurements, our improvement is built on leveraging powerful techniques such as the quantum singular value transformation (QSVT) \cite{gilyen2019qsvt, Gilyn2019thesis} and density matrix exponentiation \cite{lloyd2014QuantPrincCompAnaly}, enabling us to perform the Uhlmann transformation without state measurements.




In addition, we provide a lower bound on the query and sample complexities for the Uhlmann transformation, based on a result from the mixedness testing---the task of certifying whether a state is the completely mixed state or far from it \cite{ODonnell2021SpectrumTesting, MdW16surveyquantproptest, Wright2016howtolearnQS}.
Our result shows that, for a pair of pure states $\ket{\rho}^{\sA\sB}$ and $\ket{\sigma}^{\sA\sB}$ whose reduced states $\rho^\sA$ and $\sigma^\sA$ have the same rank, any quantum algorithm for the Uhlmann transformation necessarily requires query or sample complexity that scales as the cube root of the rank.
Together with the explicit construction of the algorithms, the lower bound allows us to characterize the algorithmic features of the Uhlmann transformation in terms of both upper and lower bounds on its computational costs.



We apply the Uhlmann transformation algorithms to fidelity estimation, the task of estimating the fidelity between two states $\rho^\sA$ and $\sigma^\sA$.
This task is crucial for information processing, as it allows us to assess the accuracy of quantum state preparation, the reliability of quantum communication, and the outputs of quantum circuits and algorithms.
We present estimation algorithms for the square root fidelity $\sqrt{\rF}(\rho^\sA, \sigma^\sA)$ based on the Uhlmann transformation and evaluate their computational costs in the three query and sample models.
Notably, our algorithm achieves a significant improvement in both query and sample complexities over the prior state-of-the-art algorithms in Refs.~\cite{gilyen2022improvedfidelity, Liu2024geometricmean}.


The key idea behind our approach is that, according to the Uhlmann's theorem in Eq.~\eqref{inteq:57}, we reduce the problem of estimating the fidelity between generally mixed states $\rho^\sA$ and $\sigma^\sA$ to estimating the fidelity between pure states.
We apply the Uhlmann transformation using our algorithms, followed by known fidelity estimation techniques for pure states \cite{Wang2024OptimalFidelityPure, wang2024sampleoptimal}, which achieve an optimal number of queries or samples.
In the mixed sample access model, this approach requires the given mixed states to be purified. To this end, we employ purification subroutines that prepare a purification from multiple samples, such as the random purification \cite{tang2025conjugatequerieshelp} or the so-called canonical purification.
Consequently, we obtain improved algorithms for square root fidelity estimation between general states.
This also serves as a good application of the recently proposed subroutine in Ref.~\cite{tang2025conjugatequerieshelp}, leading to a constructive algorithmic result.




To highlight the wide-ranging utility of our algorithms, we further explore other applications to several information-theoretic tasks. 
Based on the decoupling approach \cite{hayden2008decoupling, dupuis2010decoupling, Szehr2013decouple2design, dupuis2014one, preskill2016quantumshannon}, the Uhlmann transformation is known as an operation that can achieve theoretical limits, such as the quantum capacity \cite{hayden2008decoupling}.
In this study, we focus on two tasks: entanglement transmission \cite{Schumacher1996sendingentanglement, schumacher2001approximateerrorcorrection, hayden2007black, datta2013oneshotentassistQCcom, beigi2016decoding, khatri2024principlemodern} and quantum state merging \cite{Horodecki2005Partialquantinfo, berta2008singleshot, dupuis2014one, yamasaki2019statemergesmalldimension}.
We apply our Uhlmann transformation algorithms to accomplish these tasks in the form of quantum circuits and evaluate their computational costs.


We also apply our algorithms to the algorithmic implementation of the Petz recovery map \cite{petz1986sufficient, petz1988sufficiency}, a map which plays an important role in various fields of quantum science \cite{barnum2002reversing, ng2010simpleapproach, chen2020entanglement, lauten2022approx, gilyen2022petzmap, biswas2023noiseadapted, nakayama2023petz}.
The Petz recovery map is defined for a quantum channel and a reference state, and is known to reverse the effect of the channel with near-optimal error \cite{barnum2002reversing, beigi2016decoding}.
While quantum circuits for implementing the Petz recovery map were previously studied \cite{gilyen2022petzmap, biswas2023noiseadapted}, they relied on a crucial assumption that a Stinespring dilation unitary \cite{stinespring1955} of the channel and its inverse can be implemented. 
We provide an algorithm without this assumption---which approximately realizes the Petz recovery map using only the quantum channel itself and a reference state---by leveraging our Uhlmann transformation algorithm for the mixed sample access model.
The core subroutine that makes this possible is an algorithm that applies a Stinespring unitary of a quantum channel.
This algorithm is composed of the canonical purification algorithm and the Uhlmann transformation algorithm.




The rest of this paper is organized as follows. An overview of our main results is provided in \Cref{sec:mainresults}, and summary and outlooks are in \Cref{sec:conclusion}. 
We provide thorough preliminaries in \Cref{sec:preliminaries}.
The details of the algorithms in each computational model are explained in Secs.~\ref{sec:purif query},~\ref{sec:purif sample}, and~\ref{sec:mixed sample}, respectively.
Applications of our algorithm are presented in \Cref{sec:fidelity estimation} for fidelity estimation, in \Cref{sec:decoupling and uhlmann} for the combination with the decoupling approach, and in \Cref{sec:sample petz} for the algorithmic implementation of the Petz recovery map.
Technical materials are provided in Appendices~\ref{sec:altanative mixed sample uhlmann} to~\ref{sec:appendix error not accumulate}. 









